% !TeX root = ../nc-quadric-avatar.tex
% 2026-11-22: 幾何的avatarの四つのcanonical列
\section{反標準 avatar の canonical complexes と周期}
\label{sec:qa-self-contained-helix}

この節は前節の幾何的補題から四項の canonical complex を
構成する．これは有限 foundation と全周期の関係を明示する
ためのものであり，任意の相対基底上の AS 正則性と
内在的な心の同定を無条件に主張するものではない．

\begin{proposition}[rotated square の四つの完全列]
\label{prop:qa-core-rotated-complexes}
定理~\ref{thm:qa-core-avatar-square-rotated} の
rotated square 曲線型 avatar を固定し，
\[
\begin{gathered}
L_0=P,\quad L_1=HF^{-2},\quad
L_2=HF^{-1},\quad L_3=H^2F^{-1},\\
L_4=H^2P,\quad L_5=H^3F^{-2},\quad
L_6=H^3F^{-1}P^2,\quad L_7=H^4F^{-1}P^6 .
\end{gathered}
\]
これらの線束の間には directed foundation の評価写像と
最小関係から作られる次の四つの局所分裂完全列がある：
\begin{align}
0&\to L_0\to L_2^{\oplus2}\to L_3^{\oplus2}
     \to L_4\to0,\label{eq:qa-core-rsq-complex0}\\
0&\to L_1\to L_2^{\oplus2}\to L_3^{\oplus2}
     \to L_5\to0,\label{eq:qa-core-rsq-complex1}\\
0&\to L_2\to L_3^{\oplus4}
     \to L_4^{\oplus2}\oplus L_5^{\oplus2}
     \to L_6\to0,\label{eq:qa-core-rsq-complex2}\\
0&\to L_3\to L_4^{\oplus2}\oplus L_5^{\oplus2}
     \to L_6^{\oplus4}\to L_7\to0.
     \label{eq:qa-core-rsq-complex3}
\end{align}
任意の $q\in\mathbb Z$ に対して，$i=0,1,2,3$ とし
\[
\begin{array}{c|cccc}
i&0&1&2&3\\ \hline
G_i&0&h-2f&h-f&2h-f\\
d_i&0&0&2&6\\
\varepsilon_i&1&0&0&0
\end{array}
\]
と書くと，
\begin{equation}
L_{4q+i}=\OO_D(G_i+2qh)
           P^{\,4q(q-1)+d_iq+\varepsilon_i}
\label{eq:qa-core-rsq-periodic-L}
\end{equation}
は同じ四列の全整数平行移動を持つ．
\end{proposition}
\begin{proof}
$A=L_2L_0^{-1}=HF^{-1}P^{-1}$ は
補題~\ref{lem:qa-core-curve-vanishing} により二切断の
大域生成 pencil である．評価列
$0\to A^{-1}\to\OO_D^{\oplus2}\to A\to0$
を $L_2$ で tensor すると
\[
0\to L_0\to L_2^{\oplus2}\to L_2^2L_0^{-1}\to0 .
\]
他方 $L_3(L_2^2L_0^{-1})^{-1}=FP$ も pencil なので，
その評価列を $L_3$ で tensor した
\[
0\to L_2^2L_0^{-1}\to L_3^{\oplus2}\to L_4\to0
\]
と splice すれば第一列となる．
第二列は $L_2L_1^{-1}=F$ と
$L_3H^{-1}=HF^{-1}$ という二 pencil の評価列
\[
0\to L_1\to L_2^{\oplus2}\to H\to0,\qquad
0\to H\to L_3^{\oplus2}\to L_5\to0
\]
を splice して得られる．
ここで $HF^{-1}$ は二切断の大域生成 pencil である．

第三列には，二つの pencil $a_+=FP$，$b_+=HF^{-1}$
を用いる．比 $b_+a_+^{-1}=HF^{-2}P^{-1}$ は
補題~\ref{lem:qa-core-curve-vanishing} より acyclic，
積 $a_+b_+=HP$ は very ample である．
よって補題~\ref{lem:qa-core-two-pencils} で得る
別の $(2,2)$ 埋め込み上に，
命題~\ref{prop:qa-core-directed-foundations} の
四行四列の対角分解を適用する．
その曲線への導来制限では，
Cartier 因子の標準 Tor 項
$\operatorname{Tor}_1(\iota_*Q,\OO_D)=Q(-D)$
が左端に現れる．
$Q=HP^2$ を代入し，全体を $L_3$ で tensor すると
第三列の四項と係数 $4,2+2,1$ が得られる．
これは局所的に長さ一の Cartier 解消を制限して
得た完全列なので，単に Euler 特性の等号ではない．

第四列も，二つの pencil
$a_0=HF^{-1}P$，$b_0=FP^2$ を使う．
その比 $a_0b_0^{-1}=HF^{-2}P^{-1}$ は acyclic，
積は $HP^3$ で very ample である．
同じ四行四列分解を $Q=HP^2$ に適用し，
導来制限と双対化を行った後 $L_6$ で tensor すると
第四列を得る．
第三，第四列の中間写像は対角分解の評価から得られ，
大域切断に戻すと四本の矢と二本ずつの最小関係に一致する．
全列は線束の局所分裂完全列である．

周期については，各連続四項を共通線束で
tensor して最初の項を $P$ に正規化すると，
二つの ruling の役割を担う線束は
$F_q=FP^{2q}$ と $H_q=HP^{4q}$ となる．
$H_qF_q^{-2}=HF^{-2}$ は不変なので，これらは
同じ非自明性と消滅条件を満たす．
前節の二 pencil の再埋め込みを各 $q\in\mathbb Z$
に適用でき，上の四列を繰り返して構成できる．
線束の差分を帰納的に追うと
式~\eqref{eq:qa-core-rsq-periodic-L} を得る．
\end{proof}

\begin{corollary}[square の四つの完全列]
\label{cor:qa-core-square-complexes}
square avatar の $L_0=P,L_1=F,L_2=G,L_3=FG$ に対し
\[
(L_4,L_5,L_6,L_7)
=(F^2G^2P,\ F^3G^2P^2,\ F^2G^3P^2,\ F^3G^3P^4)
\]
とおくと，
\begin{align*}
0&\to L_0\to L_1^{\oplus2}\oplus L_2^{\oplus2}
       \to L_3^{\oplus4}\to L_4\to0,\\
0&\to L_1\to L_3^{\oplus2}\to L_4^{\oplus2}
       \to L_5\to0,\\
0&\to L_2\to L_3^{\oplus2}\to L_4^{\oplus2}
       \to L_6\to0,\\
0&\to L_3\to L_4^{\oplus4}
       \to L_5^{\oplus2}\oplus L_6^{\oplus2}
       \to L_7\to0
\end{align*}
が canonical な局所分裂完全列となる．
\end{corollary}
\begin{proof}
定理~\ref{thm:qa-core-avatar-square-rotated} により，
square の一つ後方または前方の foundation を
rotated の四項と比較する．
正確には rotated の $(L_{-1},L_0,L_1,L_2)$ を
$L_{-1}=F^{-1}P^2$ から出発して
共通線束 $FP^{-1}$ で tensor すると
$(P,F,G,FG)$ となる．
前命題の列を一つ添字移動し，
全項にこの共通線束を tensor すれば四列を得る．
この操作は曲線上の線束の標準化であって，
非可換圏に $(-)\otimes FP^{-1}$ という自己同値を
仮定したものではない．
\end{proof}

\begin{proposition}[全整数 helix の前向き Ext 消滅]
\label{prop:qa-core-forward-ext-vanishing}
上記の square および rotated square の曲線上の線束列について
\[
 \Ext_D^p(L_i,L_j)=0\qquad(p>0,\ i<j)
\]
が成り立つ．曲線が可約・非被約であっても同様である．
これは曲線上の forward acyclicity の主張であり，
この段階で例外 collection の fullness や
非可換 AS--Gorenstein 性を主張するものではない．
\end{proposition}
\begin{proof}
Gorenstein 曲線 $D$ 上で
$\Ext^p_D(L_i,L_j)=H^p(D,L_jL_i^{-1})$ であり，
$p\ge2$ は自動的に零である．
rotated の連続する四商は，$q$ 番目の区間で
\[
 C P^{-1},\quad FP^{2q},\quad HP^{4q},
       \quad FP^{2q+1},
 \qquad C=HF^{-2}.
\]
唯一の総次数零の商 $CP^{-1}$ は
avatar の除外条件 $P\ne C$ と
補題~\ref{lem:qa-core-curve-vanishing}(4) で acyclic である．
その他の商は各成分上非負次数の線束であり，
少なくとも一つ含めば全体上の次数は正である．
複数の連続商の tensor 積は，単独の
$CP^{-1}$ 以外は各部分曲線上非負次数を持ち，
$D$ 全体上正次数を持つ．
ここで $C$ 自体の成分次数は零ではない場合があるが，
$C$ の両側には $F$ または $H$ の商が隣接し，
それらとの積は各部分曲線に非負次数を持つ．
補題~\ref{lem:qa-core-curve-vanishing}の証明で得た
cluster 消滅の不等式 $(**)$ を $\ell=0$ に適用すれば
$H^1=0$ である．

square では四つの連続商は
\[
 FP^{2q-1},\quad GF^{-1},\quad FP^{2q},
       \quad FGP^{4q+1}
\]
である．$GF^{-1}$ の単独の acyclicity は
$\Sigma_0$ 上の Cartier 制限列，
または補題~\ref{lem:qa-core-curve-vanishing}(4)
による．それ以外は上と同様に
nef な正次数の商を少なくとも一つ含み，
$GF^{-1}$ を隣接商と組にすると
各部分曲線上の次数は非負になる．
同じ cluster 判定から $H^1=0$ が従う．
\end{proof}

\begin{remark}[AS 正則性への論理的な境界]
\label{rem:qa-core-helix-not-as}
前命題と系は，両端と中間項を持つ
\emph{実際に完全な}曲線上の複体を与える．
非可換な全 $\mathbb Z$-代数の AS 正則性には，
これらを canonical exceptional helix の
最小射影分解として実現し，
反対向きの $\Ext^3$ による Gorenstein 性を示す
別の導来圏比較が必要である．
この比較のためには acyclic spherical helix と
有限 directed algebra の Serre 関手を用いる．
その\emph{体上の全AS正則$\mathbb Z$-代数への橋}を
第~\ref{sec:qa-as-avatar-classification}節で
最小分解の必要条件・逆構成・相互逆性まで証明する．
ただし相対基底変更と元の O-marked heart の直接同値は
この四列だけからは結論しない．
\end{remark}
